% Original Markdown SHA-256: 2d443b8fed3a39951bedbe233458734993909828a480b78297cac94032df88ca
\chapter{Exact graphical direction modules and supported arithmetic visibility}
\label{chapter:11_graphic-visibility-filtration}
{\small\noindent\textbf{Source:} \nolinkurl{history/EP817_SUPPORTED_DEFECT_20260915/payloads/ep817/research/supported_defects/notes/graphic-visibility-filtration.md}\
\textbf{Identity:} \nolinkurl{2d443b8fed3a39951bedbe233458734993909828a480b78297cac94032df88ca}}
\medskip
15 September 2026. Research continuation with ordinary proofs and bounded exact certificates; no new Lean elaboration or priority claim. The cut-tree existence theorem used here is classical and is attributed below. The original graph, its orientations, the integral cut capacities, and the arithmetic observation remain in every comparison.

\section{1. The original obstruction and the stronger question}\label{n11-the-original-obstruction-and-the-stronger-question}

The preceding complete-graph argument uses edge-disjoint paths: r paths joining u to v give r+1 subset sums in progression, provided the observed endpoint difference t\_v-t\_u is nonzero. This contribution asks what happens when that difference is a supported zero and how to account for \emph{every} progression direction of an arbitrary finite graph, rather than only a chosen path family.

Let G=(V,E) be a finite loopless undirected multigraph. Parallel edges are separate source labels. Choose and retain an orientation of each edge, with incidence

\[
\begin{adjustbox}{max width=.92\linewidth}
$\displaystyle \partial:\mathbb Z^E\to\mathbb Z^V.$
\end{adjustbox}
\]

Every binary choice reverses the corresponding edge of a reference orientation. Let Sigma(G) be the set of its actual orientation-score vectors; their coordinates count wins, and their sum on each connected component is its number of edges. Let mu\_G(s) be the number of reference reversal masks producing score s. Those multiplicities are retained.

For integer vertex marks T=(t\_v), define

\[
\begin{adjustbox}{max width=.92\linewidth}
$\displaystyle L_T(s)=\sum_v t_v s_v.$
\end{adjustbox}
\]

If every chosen edge orientation goes from smaller to larger mark, and all edge distances are distinct positive integers, the original subset-sum image is a reflection/translation of L\_T(Sigma(G)). Specifically, with reference score rho,

\[
\begin{adjustbox}{max width=.92\linewidth}
$\displaystyle \Phi_A(\epsilon)=L_T(\rho)-L_T(\rho-\partial\epsilon).$
\end{adjustbox}
\tag{1.1}
\]

Equal marks at nonadjacent vertices are allowed by this identity. Cardinality claims about the EP817 generators still require all actual edge distances to be distinct; no multigraph edge count is substituted for that condition.

Write lambda\_G(u,v) for the maximum number of edge-disjoint u-v paths, equal to the minimum cardinality of an original separating edge cut. Put lambda\_max(G)=max lambda\_G(u,v), with value zero for an edgeless graph.

For q\textgreater=1, define the integral module generated by actual source progression directions:

\[
\begin{adjustbox}{max width=.92\linewidth}
$\displaystyle \mathcal D_q(G)=\left\langle z\in\mathbb Z^V:
 \exists x\in\Sigma(G),\quad x,x+z,\ldots,x+qz\in\Sigma(G)
 \right\rangle_{\mathbb Z}.$
\end{adjustbox}
\tag{1.2}
\]

The bracket is explicitly the integral span; the actual feasible direction set is retained before taking that span.

\section{2. The classical cut tree and its exact source-coordinate map}\label{n11-the-classical-cut-tree-and-its-exact-source-coordinate-map}

For each nontrivial connected component, the Gomory-Hu theorem supplies a weighted tree on the \emph{same original vertices}. Removing a tree edge e gives an actual vertex partition U\_e and its complement. The weight is

\[
\begin{adjustbox}{max width=.92\linewidth}
$\displaystyle c_e=|\delta_G(U_e)|,$
\end{adjustbox}
\]

and for all vertices u,v,

\[
\begin{adjustbox}{max width=.92\linewidth}
$\displaystyle \lambda_G(u,v)=\min_{e\in\text{tree path }u\leadsto v}c_e.$
\end{adjustbox}
\tag{2.1}
\]

Each fundamental tree cut is an actual minimum cut for the endpoints of that edge. This existence/representation theorem is the imported classical result (Gomory \& Hu, 1961; Gusfield, 1990). No new algorithmic complexity claim is made. The finite checker supplies and independently verifies the entire cut tree on each enumerated graph.

Fix a root and choose U\_e as the subtree away from that root. For a score vector s, retain the cut coordinates s(U\_e), together with the fixed component score total. These coordinates have an explicit integral inverse: at a nonroot vertex subtract the cut sums of its child subtrees from the sum of its own subtree; at the root subtract all child subtree sums from the retained total. Thus the cut-coordinate map is an integral affine isomorphism on the component-total hyperplane, not an untracked reduction of dimension.

Every orientation satisfies

\[
\begin{adjustbox}{max width=.92\linewidth}
$\displaystyle |E(G[U_e])|\le s(U_e)
 \le |E(G[U_e])|+c_e.$
\end{adjustbox}
\tag{2.2}
\]

Internal edges contribute exactly one win to U\_e; crossing edges contribute zero or one. Hence the displayed interval is the exact original range, of width c\_e.

\section{3. The full direction filtration}\label{n11-the-full-direction-filtration}

Delete from the cut forest every tree edge with c\_e\textless q. Let P\_q be the resulting partition of V. By (2.1), two vertices belong to the same class exactly when lambda\_G(u,v)\textgreater=q. The partition therefore does not depend on which valid cut tree was chosen.

Define

\[
\begin{adjustbox}{max width=.92\linewidth}
$\displaystyle K_q(G)=\left\{z\in\mathbb Z^V:
             \sum_{v\in C}z_v=0\text{ for each }C\in P_q\right\}.$
\end{adjustbox}
\tag{3.1}
\]

It is the direct sum of the zero-sum lattices of the actual classes, of rank \textbar V\textbar-\textbar P\_q\textbar.

\textbf{Theorem 3.1 (exact direction module).} For every finite G and every q\textgreater=1,

\[
\begin{adjustbox}{max width=.92\linewidth}
$\displaystyle \boxed{\mathcal D_q(G)=K_q(G).}$
\end{adjustbox}
\tag{3.2}
\]

Every root generator e\_v-e\_u in one class of P\_q has an explicitly constructible (q+1)-term progression in the original orientation source.

\textbf{Proof of the upper inclusion.} Take an actual source progression x,x+z,...,x+qz. For a cut-tree edge of weight c\_e\textless q, (2.2) gives

\[
\begin{adjustbox}{max width=.92\linewidth}
$\displaystyle q|z(U_e)|\le c_e<q.$
\end{adjustbox}
\]

The integer z(U\_e) is therefore zero. Component totals also give zero total charge on every original component. Contract the high-weight tree edges. The remaining graph is a tree on the classes P\_q, and every one of its fundamental cut charges is zero. The explicit leaf/subtree inverse of Section 2 now gives zero charge on each class. Thus z belongs to K\_q, and the same holds for its integral span.

\textbf{Proof of the lower inclusion.} If u,v belong to one class, unit-capacity max-flow supplies q edge-disjoint u-v paths. It can be decomposed into q integral path chains c\_h with boundary e\_v-e\_u; a standard augmenting-path construction supplies actual paths, rather than fractional capacities. Write each \(c_h=c_h^+-c_h^-\) using its positive and negative edge supports relative to the retained reference orientation.

Choose c\_h\^{}+ on all paths, then replace them one at a time by c\_h\^{}-. Disjointness of the paths makes every union an actual binary edge choice. Its boundary decreases by e\_v-e\_u at every step, so its score increases by e\_v-e\_u. This constructs the promised (q+1)-term progression. The vectors e\_v-e\_u in each class generate its zero-sum lattice, proving the lower inclusion and the theorem. QED.

The distinction between directions and their integral span remains explicit. The proof supplies actual progressions for the class-root generators; it does not claim that every integral combination is itself a feasible direction.

\textbf{Corollary 3.2.} The longest progression in Sigma(G), allowing a singleton when no nonconstant progression exists, has length

\[
\begin{adjustbox}{max width=.92\linewidth}
$\displaystyle \boxed{1+\lambda_{\max}(G).}$
\end{adjustbox}
\tag{3.3}
\]

The lower witness is the path construction. Above that length, P\_q is discrete, K\_q=0 and (3.2) excludes a nonzero direction. The entire unbounded q-filtration changes at only the finitely many cut-tree weights; the original q label is retained on any zero or constant plateau.

\section{4. Exact supported visibility under the arithmetic observation}\label{n11-exact-supported-visibility-under-the-arithmetic-observation}

For each q retain

\[
\begin{adjustbox}{max width=.92\linewidth}
$\displaystyle N_{q,T}=K_q\cap\ker L_T,\qquad
 I_{q,T}=L_T(K_q)\subseteq\mathbb Z.$
\end{adjustbox}
\]

The actual image ideal is

\[
\begin{adjustbox}{max width=.92\linewidth}
$\displaystyle \boxed{I_{q,T}=a_q\mathbb Z,\qquad
 a_q=\gcd\{|t_v-t_u|:u,v\text{ in the same class of }P_q\},}$
\end{adjustbox}
\tag{4.1}
\]

with gcd of the empty/all-zero family equal to zero. This follows from the root generators in Theorem 3.1. The full original sequence is

\[
\begin{adjustbox}{max width=.92\linewidth}
$\displaystyle \boxed{0\to N_{q,T}\to K_q\xrightarrow{L_T}a_q\mathbb Z\to0.}$
\end{adjustbox}
\tag{4.2}
\]

The terminal ideal has its original inclusion in Z. It is not divided by a\_q and then silently treated as a fixed metric copy of Z. The gcd describes the image of the integral span; individual witnessed progressions retain their actual endpoint differences.

\textbf{Theorem 4.1 (visibility).} A (q+1)-term progression in the original score source has nonconstant numerical image under L\_T if and only if a\_q\textgreater0. Equivalently, the largest visible source progression has length

\[
\begin{adjustbox}{max width=.92\linewidth}
$\displaystyle \boxed{1+\max_{u,v:t_u\ne t_v}\lambda_G(u,v),}$
\end{adjustbox}
\tag{4.3}
\]

with maximum zero when there is no distinct observed pair.

\textbf{Proof.} A nonconstant observed source progression has a direction z with L\_Tz!=0. Theorem 3.1 puts z in K\_q, so its image ideal is nonzero. Conversely, a\_q\textgreater0 gives u,v in a common class with unequal marks. The root progression constructed in Theorem 3.1 has observed difference t\_v-t\_u!=0. The equivalence and the maximum formula follow. QED.

When a\_q=0, equation (4.2) says N\_(q,T)=K\_q; it does not say that K\_q has disappeared. On an active support its observed coefficient is the supported zero. This is a direct instance of the Zeta source\textquotesingle s receiving-boundary and observation-kernel formula (D7)--(D8), with all source directions still present.

For q increasing, K\_(q+1) is a sublattice of K\_q. Every inclusion commutes with L\_T, giving morphisms of the original kernel-image sequences. The entire infinite index family has a literal SplitZero realization: take

\[
\begin{adjustbox}{max width=.92\linewidth}
$\displaystyle \mathcal L=(\{1,2,\ldots\}\cup\{\infty\},\ \text{reverse usual order}),
 \quad\bot=\infty,\quad q\vee p=\min(q,p),
 \quad K_\infty=0.$
\end{adjustbox}
\tag{4.4}
\]

For q\textgreater=p, the transport K\_q-\textgreater K\_p is the original lattice inclusion. Use the same supports for N\_(q,T) and I\_(q,T), and their original inclusions. Reconstructing these diagrams with the source\textquotesingle s G(Z) functor retains every finite q label. In particular, K\_q=0 at a finite q is its supported zero fibre; it is not the bottom at infinity.

There is also a finite support-changing map. Let P be the finite chain of actually occurring cut partitions, with a \emph{separate} external bottom adjoined. Send infinity to that bottom and a finite q to the active partition P\_q. This preserves bottom and joins, since P\_min(q,p) is the coarser of P\_q and P\_p. Its coefficient map is the literal identity K\_q=K\_(P\_q); on reconstructed totals it is

\[
\begin{adjustbox}{max width=.92\linewidth}
$\displaystyle (q,z)\longmapsto(P_q,z).$
\end{adjustbox}
\tag{4.5}
\]

Two images agree exactly when their cut partitions agree and their coefficient vectors agree. Thus the map\textquotesingle s support coalescences are specified, while the original unbounded q values remain on the source. The active discrete partition at q\textgreater lambda\_max is distinct from the separately adjoined bottom. No finite supported zero is sent to external absence. The observation and all kernel-image maps commute with (4.5). This is finite control of the original infinite filtration with its comparison map and support fibres exposed; no analytic theta estimate is used.

\section{5. The positive arithmetic consequences and the retained extra kernel}\label{n11-the-positive-arithmetic-consequences-and-the-retained-extra-kernel}

If all vertex marks are globally distinct and the observed subset-sum image is k-AP-free, (4.3) forces

\[
\begin{adjustbox}{max width=.92\linewidth}
$\displaystyle \lambda_{\max}(G)\le k-2.$
\end{adjustbox}
\tag{5.1}
\]

For a simple graph let r=\textbar V\textbar-c(G), the rank of its incidence lattice. Map each original edge to its signed path in the cut forest. A column has at least one nonzero entry. At most r columns have exactly one entry, since a simple original edge can connect the endpoints of a given tree edge at most once. Counting every entry in both orders gives

\[
\begin{adjustbox}{max width=.92\linewidth}
$\displaystyle 2|E|-r\le\sum_{a\in E}|\text{tree path of }a|
 =\sum_{e\in\text{cut forest}}c_e
 \le\lambda_{\max}(G)r.$
\end{adjustbox}
\]

Consequently

\[
\begin{adjustbox}{max width=.92\linewidth}
$\displaystyle \boxed{2|E|\le(\lambda_{\max}(G)+1)r,\qquad
       2|E|\le(k-1)r\text{ under (5.1)}.}$
\end{adjustbox}
\tag{5.2}
\]

This workbench deduction uses the classical cut-tree representation. It strengthens the earlier elementary graphical edge-count consequence, but no priority claim is made for the graph inequality itself. Parallel original edges would repeat unit columns; the simple-graph condition is therefore retained in this count.

A scalar projection can create progressions whose chosen source lifts are not progressions. The comparison is the following exact pullback, for q\textgreater=2:

\[
\begin{adjustbox}{max width=.92\linewidth}
$\displaystyle \mathcal F_T=
 \{(s_0,...,s_q)\in\Sigma(G)^{q+1}:
             \Delta(s_0,...,s_q)\in(\ker L_T)^{q-1}\}.$
\end{adjustbox}
\tag{5.3}
\]

The actual vector-progression set is the subfibre Delta=0. Evaluation maps F\_T onto the complete scalar progression set; each numerical tuple has its actual score-lift fibre, and each score in turn has its original orientation masks. The visibility theorem calculates the Delta=0 part with nonzero observed difference. The remaining feasible fibre in (5.3) is retained rather than being identified with that part.

For example a two-edge star with marks 0,1,2 has source maximum progression length two, whereas its scalar subset sums are \{0,1,2,3\}. The source tuples with scalar three-APs are accounted for by the nonzero receiving defects in (5.3). This explicit comparison prevents using (5.1) as an unproved converse.

\subsection{5.1 The observation-faithful quotient, with all original edges retained}\label{n11-the-observation-faithful-quotient-with-all-original-edges-retained}

There is a further exact comparison when the graph actually represents an EP817 set. Assume every edge has nonzero mark difference and all absolute edge differences are pairwise distinct. Identify original vertices exactly when their arithmetic marks agree. Let p:V-\textgreater Vbar be this quotient, and retain each original edge as a separately labelled edge joining its two image vertices.

No edge becomes a loop, because an edge has nonzero mark difference. No two retained edges become parallel: that would give the same unordered pair of endpoint marks and therefore the same absolute difference, contrary to the stated distinctness. Thus Gbar is a simple graph with the original edge set, and its vertex marks are globally distinct. Its generator set is literally the same set of positive integers, not just an image with the same cardinality.

The chain square is

\[
\begin{adjustbox}{max width=.92\linewidth}
$\displaystyle \begin{array}{ccc}
\mathbb Z^E&\xrightarrow{\partial_G}&\mathbb Z^V\\
\mathrm{id}\downarrow&&\downarrow p_*\\
\mathbb Z^E&\xrightarrow{\partial_{\bar G}}&\mathbb Z^{\bar V},
\end{array}
\qquad L_T=L_{\bar T}p_*.$
\end{adjustbox}
\tag{5.4}
\]

For every orientation mask, its quotient score is exactly p\_*s: the wins at vertices in one observed class are added, with every incident edge still present. All quotient orientations lift by the identity on the edge-mask set. Consequently

\[
\begin{adjustbox}{max width=.92\linewidth}
$\displaystyle p_*:\Sigma(G)\twoheadrightarrow\Sigma(\bar G),\qquad
 \mu_{\bar G}(y)=\sum_{s:p_*s=y}\mu_G(s).$
\end{adjustbox}
\tag{5.5}
\]

The cycles created by this vertex identification are retained by the exact sequence

\[
\begin{adjustbox}{max width=.92\linewidth}
$\displaystyle \boxed{0\to H_1(G;\mathbb Z)\to H_1(\bar G;\mathbb Z)
 \xrightarrow{\partial_G}
       \operatorname{im}\partial_G\cap\ker p_*\to0.}$
\end{adjustbox}
\tag{5.6}
\]

Here graphs are their original one-dimensional chain complexes, so H\_1 is the literal incidence kernel. A quotient cycle has original boundary in the displayed intersection. The kernel of that boundary is H\_1(G). Every z in the intersection has an original edge lift c with partial\_G c=z, and p\_*z=0 makes c a quotient cycle; this proves surjectivity over the integers, not only equality of rational dimensions.

If r\_G=\textbar V\textbar-c(G) and r\_bar=\textbar Vbar\textbar-c(Gbar), the last free lattice has rank r\_G-r\_bar. Its arithmetic observation vanishes, since L\_T=L\_barT p\_*. These are present original zero actions, not deleted edge choices. For the Theta\_r family, (5.6) adds exactly one cycle direction when the two equally marked terminals are identified; the quotient is the original bouquet of r triangles.

The word metric also commutes with this comparison. On the score values use the original quotient Gram diag(1/mu\_G(s)); then the least-norm section of (5.5) is

\[
\begin{adjustbox}{max width=.92\linewidth}
$\displaystyle e_y\longmapsto\sum_{p_*s=y}\frac{\mu_G(s)}{\mu_{\bar G}(y)}e_s.$
\end{adjustbox}
\tag{5.7}
\]

Its squared norm is 1/mu\_Gbar(y). Thus the minimum-norm quotient is precisely the one obtained from the same original edge masks, with no unweighted-score metric silently substituted.

The visibility theorem now applies to Gbar with its globally distinct marks. In particular every k-admissible edge-distance set with this presentation obeys

\[
\begin{adjustbox}{max width=.92\linewidth}
$\displaystyle \boxed{2|E|\le(k-1)\bigl(|\bar V|-c(\bar G)\bigr).}$
\end{adjustbox}
\tag{5.8}
\]

It counts the actual distinct observed vertex marks after retaining the full relative sequence (5.6). The earlier bound used globally distinct marks at the source; (5.8) covers coincident nonadjacent marks through their explicit quotient. A general positive generator set still has a star presentation with \textbar Vbar\textbar-c=n, so this necessary inequality alone does not yield the sought sharp global exponential rate.

\section{6. Application to an unbounded supported-zero obstruction}\label{n11-application-to-an-unbounded-supported-zero-obstruction}

Use the graph Theta\_r of the companion note: r internally disjoint paths of length three joining u to v, with the two internal vertices on path j marked L19\^{}j and 8L19\^{}j. Initially t\_u=t\_v=0. The actual edge weights are L times the original base-19 triples, all positive and distinct.

For r\textgreater=2, every vertex pair has at least two edge-disjoint paths. Pairs containing an internal vertex have at most two because that vertex has degree two. The terminal pair has exactly r disjoint paths. Thus for r\textgreater=3:

\[
\begin{adjustbox}{max width=.92\linewidth}
$\displaystyle P_1=P_2=\{V\},$
\end{adjustbox}
\]

and for 3\textless=q\textless=r the only nonsingleton class of P\_q is \{u,v\}. Above r every class is a singleton. The source module in that middle interval is exactly

\[
\begin{adjustbox}{max width=.92\linewidth}
$\displaystyle K_q=\mathbb Z(e_v-e_u).$
\end{adjustbox}
\tag{6.1}
\]

It is nonzero at arbitrarily large q as r increases. Its observation is nevertheless zero because t\_v=t\_u. Theorem 4.1 proves that the longest numerically visible \emph{source} progression has length three, uniformly over all r. The original mod-19 lift independently excludes any additional scalar four-AP in (5.3), so the actual arithmetic image is four-AP-free at every r.

Now change t\_v to -t. On the entire middle filtration,

\[
\begin{adjustbox}{max width=.92\linewidth}
$\displaystyle L_T(e_v-e_u)=-t,\quad I_{q,T}=t\mathbb Z.$
\end{adjustbox}
\tag{6.2}
\]

For t!=0 all these retained directions become visible; their path witnesses produce a progression with r+1 terms. With L\textgreater2r\textbar t\textbar{} and t\textgreater0, the companion deformation theorem proves the exact maximal length, retaining every original scalar value and representation fibre.

This controls the infinite family rather than inferring from a single rank. The graphical direction module, its supported zero observation, the relative homology under terminal identification, and the first-order specialization map are all explicit. The terminal identification leaves the literal generator set unchanged. Conversely, changing its arithmetic mark changes the scalar operator, and (6.2) records exactly which long directions become visible.

\section{7. Full finite replay}\label{n11-full-finite-replay}

The proof does not extrapolate from computations. The finite suite independently calibrates its integer constructions on every labelled simple graph on 1..5 vertices: 1,099 graphs. It rebuilds all 40,398 orientation-score values across these instances, examines 1,391,037 endpoint pairs, and supplies 1,152 valid cut trees for nontrivial components. Every fundamental cut is checked in the original graph and against all required pair minima.

The producer also constructs 15,147 root-direction witnesses, verifies 8,493 original mark observations, and checks 23,118 exact image ideals. For up to four vertices the observation domain is every mark list in \{0,1,2\}\^{}V. At five vertices the three stated lists are all zero, (0,1,2,3,4), and (0,0,1,2,3). The complete records retain this domain, including the globally unequal and supported-zero cases. The independent auditor reconstructs source orientations by direct binary enumeration, uses a separately written cut calculation, and checks all these observation records. It does not import the producer.

A separate observation-quotient replay checks 1,211 admitted marked graphs among 2,197 literal candidates (all labelled graphs through five vertices, with the two stated power-mark families), 60,585 original orientation masks, and the relative incidence sequence for Theta\_r through r=12. Full theta word fibres are enumerated through r=4. The quotient producer and its separate auditor retain original edge labels, positive distinct weights, all score-fibre masses and their quotient metrics.

The complete compressed records for the first suite are \nolinkurl{certificates/supported_defect_proof.json.gz}. False image counts, missing supports, altered cut capacities and zeroed actual witnesses are rejected. The exact source identities and replay outputs are separately hashed. The full generality of Theorems 3.1 and 4.1 comes from the cut-tree and path proofs above.

\section{References and source roles}\label{n11-references-and-source-roles}

Gomory, R. E., \& Hu, T. C. (1961). Multi-terminal network flows. \emph{Journal of the Society for Industrial and Applied Mathematics, 9}(4), 551--570. \url{https://doi.org/10.1137/0109047} . Role: existence and exact fundamental-cut representation of the cut-equivalent tree, not a source of an EP817 numerical bound.

Gusfield, D. (1990). Very simple methods for all pairs network flow analysis. \emph{SIAM Journal on Computing, 19}(1), 143--155. \url{https://doi.org/10.1137/0219009} . Role: classical cut-tree construction and interpretation. No implementation from this reference is imported into the finite checker.

Hartmann, T., \& Wagner, D. (2013). \emph{Dynamic Gomory-Hu tree construction---fast and simple} (arXiv:1310.0178) {[}Preprint{]}. \url{https://arxiv.org/abs/1310.0178} . Role: independently inspected primary-source description of the cut tree as a minimum-cut basis.

KokunoYumeto. (2026). \emph{Reconstruction with changes of support index} {[}TeX source, Zeta main 58626a62cd5648e8ae7450fb54d4a4aab2330981{]}. \nolinkurl{workbenches/splitzero-tandem/tex/support_diagrams.tex}, (D6)--(D8). Role: original supported reconstruction and quotient/observation-kernel comparison.

The Clankers. (2026, September 15). \emph{Graphical obstruction propagation and a stronger six-term construction} {[}Predecessor workbench contribution{]}. Role: disjoint path-to-subset witness map and complete-graph context; retained with original attribution and evidence scope.
